\begin{mistake}{2026-07-30}{01}

\begin{problemcontent}
设 \(f(x)\) 满足 \(f'(x) + f(x) = e^{-x}\)，\(f(0) = 0\)。
(I) 求 \(y = f(x)\) 的极值与拐点；
(II) 当 \(0 < x_1 < 1 < x_2\) 时，有 \(f(x_1) = f(x_2)\)，证明：\(x_1 + x_2 > 2\)。
\end{problemcontent}

\begin{solutioncontent}
\textbf{（I）} 
原方程 \(f'(x) + f(x) = e^{-x}\) 两边同乘积分因子 \(e^x\)，得 \((e^x f(x))' = 1\)。
两边积分得 \(e^x f(x) = x + C\)。代入初值 \(f(0) = 0\) 得 \(C = 0\)，故 \(f(x) = x e^{-x}\)。

一阶导数 \(f'(x) = (1-x)e^{-x}\)。
令 \(f'(x) = 0\) 得驻点 \(x = 1\)。
当 \(x < 1\) 时 \(f'(x) > 0\)；当 \(x > 1\) 时 \(f'(x) < 0\)。
故 \(f(x)\) 在 \(x=1\) 处取得极大值 \(f(1) = \frac{1}{e}\)，无极小值。

二阶导数 \(f''(x) = (x-2)e^{-x}\)。
令 \(f''(x) = 0\) 得 \(x = 2\)。
当 \(x < 2\) 时 \(f''(x) < 0\)；当 \(x > 2\) 时 \(f''(x) > 0\)。二阶导数在 \(x=2\) 两侧变号。
故拐点为 \((2, \frac{2}{e^2})\)。

\textbf{（II）} 
由（I）知，\(f(x)\) 在 \((0, 1)\) 单调递增，在 \((1, +\infty)\) 单调递减。
要证 \(x_1 + x_2 > 2\)，即证 \(x_2 > 2 - x_1\)。
由于 \(x_1 \in (0, 1)\)，则 \(2 - x_1 \in (1, 2)\)。
又 \(x_2 > 1\)，且 \(f(x)\) 在 \((1, +\infty)\) 上单减，故只需证 \(f(x_2) < f(2-x_1)\)。
因 \(f(x_1) = f(x_2)\)，原不等式转化为证明当 \(x \in (0, 1)\) 时，恒有 \(f(x) < f(2-x)\)。

构造辅助函数 \(g(x) = f(x) - f(2-x)\)，\(x \in (0, 1]\)。
求导：
\[
\begin{aligned}
g'(x) &= f'(x) + f'(2-x) \\
&= (1-x)e^{-x} + [1-(2-x)]e^{-(2-x)} \\
&= (1-x)(e^{-x} - e^{x-2})
\end{aligned}
\]
当 \(x \in (0, 1)\) 时，\(1-x > 0\)，且 \(-x - (x-2) = 2(1-x) > 0 \implies e^{-x} > e^{x-2}\)，故 \(g'(x) > 0\)。
因此 \(g(x)\) 在 \((0, 1]\) 上单调递增，从而 \(g(x_1) < g(1) = 0\)，即 \(f(x_1) < f(2-x_1)\)。
代入 \(f(x_1) = f(x_2)\) 并利用 \((1, +\infty)\) 上的单减性，得 \(x_2 > 2-x_1\)，即 \(x_1+x_2 > 2\) 得证。
\end{solutioncontent}

\mistakeknowledge{
\begin{itemize}
  \item \textbf{核心公式/定理}：一阶线性非齐次方程凑微分法；极值点偏移问题套路。
  \item \textbf{易错点}：构造辅助函数 \(g(x) = f(x) - f(2a-x)\) 求导时，易漏掉复合函数内层求导导致的符号变化：\([f(2a-x)]' = -f'(2a-x)\)。
  \item \textbf{关联知识}：利用函数的单调性，将函数值不等式反向转化为自变量不等式。
\end{itemize}}

\end{mistake}
